\chapter{调用 Python 推理服务：POST 与 JSON 契约}

从这一关开始迁移 ML 预测。模型是原项目训练好的两个随机森林（各约 1GB），Java 读不了 Python 的 pickle 文件，所以系统拆成两个程序：

\begin{plaincode}
浏览器 --GET /api/v1/stations/42/forecast--> Spring Boot（你写的）
                                                 |
                                                 | POST /predict（内部调用）
                                                 v
                                     Python 推理服务 ml-service/（已经写好）
                                                 |
                                                 v
                                     bike_model.pkl 等 6 个模型文件
\end{plaincode}

Python 服务不用你写。它就是把原项目的 \code{forecast.py} 包成一个 FastAPI 接口，加载的还是原来那几个模型文件，预测结果和原项目完全一样（我用同样的输入逐个对比过）。

\partinfo{%
  \code{forecast/WeatherInput.java}、\code{ForecastRequest.java}、\code{HourlyPrediction.java}、\code{ForecastResponse.java}：7.1、7.3\\
  \code{forecast/ForecastClient.java}：7.2\\
  {\small（路径相对于 \code{src/main/java/ie/dublinbikes/}）}%
}{\code{./mvnw test -Dtest='Part7*'}}

\textbf{已经写好、不要改的类：}\code{ForecastProperties}、\code{ForecastPropertiesConfig}、\code{ForecastUnavailableException}、\code{StationForecast}、\code{ForecastPointResponse}、\code{StationForecastResponse}。

\begin{note}[Part 7--9 做完之前 FinalApplicationTest 会失败]
\code{FinalApplicationTest} 启动整个应用。新加的几个类还没有 \code{@Service}、\code{@Component}，所以它会报找不到 Bean，这是正常的。做完 Part 9 它会重新变绿。Part 1--6 的测试不受影响。
\end{note}

\section{知识点}

\subsection{和 Python 服务约定的 JSON（契约）}

两个服务之间靠 JSON 通信，字段名、类型都要对得上，这叫\textbf{契约}。Python 习惯用 snake\_case：

\begin{jsoncode}
请求  POST /predict
{
  "station_id": 42,
  "capacity": 30,
  "times": ["2026-10-28T18:00:00", "2026-10-28T19:00:00"],
  "weather": [ { "avg_temp": 9.0, "avg_humidity": 82.0, "avg_pressure": 1010.0, "is_raining": 1 }, ... ]
}

响应
{
  "station_id": 42,
  "predictions": [
    { "time": "2026-10-28T18:00:00", "bikes": 12.4, "docks": 17.6,
      "avg_temp": 8.0, "avg_humidity": 83.5, "avg_pressure": 1014.5, "is_raining": 0 }
  ]
}
\end{jsoncode}

\code{weather} 是可选的：不传的话，Python 服务会用自带的天气子模型按“小时 + 月份”推断天气。这一关先不传，Part 13 接入真实天气以后再传。

\subsection{整个 record 统一改命名：@JsonNaming}

Part 3 里你给每个字段单独写了 \code{@JsonProperty}。字段一多就很啰嗦。可以在 record 上加一个注解，让\textbf{所有字段}都按同一种规则转换：

\begin{java}
@JsonNaming(PropertyNamingStrategies.SnakeCaseStrategy.class)
public record Book(String bookTitle, int pageCount) {}
// 序列化成 {"book_title": "...", "page_count": 300}，反序列化时也按这个规则对应
\end{java}

\begin{note}[包名和 Part 3 不一样]
Jackson 3 里，\code{@JsonNaming} 和 \code{PropertyNamingStrategies} 在 \code{tools.jackson.databind.annotation} 和 \code{tools.jackson.databind} 包里；而 \code{@JsonProperty}、\code{@JsonInclude} 这类“基础注解”还在 \code{com.fasterxml.jackson.annotation}。IntelliJ 自动导入时注意选对。
\end{note}

\subsection{null 字段不出现在 JSON 里：@JsonInclude}

默认情况下，值为 null 的字段会被写成 \code{"weather": null}。Python 那边对 \code{null} 和“没有这个字段”处理是一样的，但让请求更干净、也更符合“可选字段”的语义，可以在 record 上加：

\begin{java}
@JsonInclude(JsonInclude.Include.NON_NULL)     // 值为 null 的字段不写进 JSON
\end{java}

\subsection{时间：LocalDateTime 和 ISO 字符串}

\code{LocalDateTime} 会自动转成 \code{"2026-10-28T18:00:00"} 这种 ISO 格式的字符串，反过来也能自动解析。Python 的 \code{datetime} 用的也是这种格式，所以两边不需要额外处理。

\code{LocalDateTime} 没有时区，表示“都柏林当地的几点”。模型是用当地时间的“小时”训练的，所以这里用它正合适。

\subsection{RestClient 发 POST 请求}

Part 3 发的是 GET。POST 多了“请求体”：

\begin{java}
Receipt receipt = restClient.post()
        .uri("/orders")
        .contentType(MediaType.APPLICATION_JSON)    // 告诉对方：请求体是 JSON
        .body(order)                                // 一个 record，Jackson 自动转成 JSON
        .retrieve()
        .onStatus(HttpStatusCode::isError, (req, res) -> { throw ...; })
        .body(Receipt.class);                       // 响应体转成 record
\end{java}

\subsection{超时}

推理服务算得慢、或者卡住时，不能让用户一直等。\code{application.yml} 里已经设好了所有 RestClient 共用的超时：

\begin{yamlcode}
spring:
  http:
    clients:
      connect-timeout: 2s     # 2 秒连不上就放弃
      read-timeout: 10s       # 连上以后，10 秒还没有响应就放弃
\end{yamlcode}

超时以后 RestClient 抛的是网络错误（\code{ResourceAccessException}），你在 7.2 里会把它转成 \code{ForecastUnavailableException}。

\subsection{带状态码的异常}

\code{ForecastUnavailableException}（已经写好）比 Part 3 的 \code{JcDecauxException} 多记了一个 \code{status}：对方返回的 HTTP 状态码，网络错误时是 0。它的 \code{isRetryable()} 根据状态码判断“值不值得重试”，Part 8 会用到。

\subsection{测试是怎么验证的}

和 Part 3 一样用 \code{MockRestServiceServer} 冒充 Python 服务。这次它还会检查\textbf{你发出的请求体}：

\begin{java}
server.expect(requestTo("http://forecast.test/predict"))
      .andExpect(jsonPath("$.station_id").value(42))       // 请求体里必须有 station_id = 42
      .andExpect(jsonPath("$.weather").doesNotExist())     // 而且没有 weather 字段
      .andRespond(withSuccess(...));
\end{java}

\section{题目}

\begin{task}{7.1\quad JSON 字段统一用 snake\_case}
\code{WeatherInput}、\code{ForecastRequest}、\code{HourlyPrediction}、\code{ForecastResponse} 四个 record 的字段已经写好了。给它们加上注解，让所有字段按 snake\_case 序列化和反序列化。

\begin{example}
\begin{plaincode}
new WeatherInput(9.0, 82.0, 1010.0, 1)
  -> {"avg_temp":9.0,"avg_humidity":82.0,"avg_pressure":1010.0,"is_raining":1}

{"time":"2026-10-28T18:00:00","bikes":12.4,...,"avg_temp":8.0,...,"is_raining":0}
  -> HourlyPrediction[time=2026-10-28T18:00, bikes=12.4, ..., avgTemp=8.0, ..., isRaining=0]
\end{plaincode}
\end{example}
\end{task}

\begin{task}{7.2\quad ForecastClient}
\textbf{a)} 交给 Spring 管理；构造函数里用注入的 \code{RestClient.Builder} 创建 RestClient，\code{baseUrl} 用 \code{properties.baseUrl()}。

\textbf{b)} \code{predict(request)}：\code{POST {baseUrl}/predict}，请求体是 \code{request}，响应转成 \code{ForecastResponse}，返回它的 \code{predictions()}。

\textbf{c)} 错误处理：
\begin{itemize}
  \item 对方返回 4xx / 5xx：抛 \code{new ForecastUnavailableException(消息, 状态码)}，消息里要有状态码；
  \item 网络错误：抛 \code{new ForecastUnavailableException(消息, 原始异常)}。
\end{itemize}

\begin{example}
\begin{plaincode}
对方返回 503  ->  status = 503, isRetryable() = true
对方返回 422  ->  status = 422, isRetryable() = false
连接被拒绝    ->  status = 0,   isRetryable() = true，cause 不是 null
\end{plaincode}
\end{example}
\end{task}

\begin{task}{7.3\quad weather 可选}
\code{ForecastRequest} 的 \code{weather} 为 null 时，JSON 里不要出现 \code{weather} 字段。

\begin{example}
\begin{plaincode}
new ForecastRequest(42, 30, times, null)
  -> {"station_id":42,"capacity":30,"times":[...]}          // 没有 "weather"
\end{plaincode}
\end{example}
\end{task}

\section{自己启动 Python 服务看看（可选）}

测试不需要 Python 服务。想亲眼看到真实的预测，在仓库根目录：

\begin{shell}
cd ml-service
.venv/bin/uvicorn app:app --port 8001          # 加载模型要 10 秒左右，看到 startup complete 就好了
\end{shell}

另开一个终端：

\begin{shell}
curl -X POST localhost:8001/predict -H 'Content-Type: application/json' \
     -d '{"station_id":42,"capacity":30,"times":["2026-10-28T18:00:00"]}'
\end{shell}

\code{.venv} 和模型文件我已经准备好了，都不进 git。换一台电脑时按 \code{ml-service/README.md} 重新准备。

\begin{think}
\begin{enumerate}
  \item 两个服务之间的 JSON 契约要改字段名时，怎么改才不会让正在运行的另一方出错？（提示：先兼容新旧两种，再删旧的）
  \item 为什么模型不在 Java 里直接加载，而要单独起一个 Python 服务？这样做有什么代价？
  \item 连接超时和读取超时有什么区别？各应该设多长？
  \item \code{@JsonNaming} 和逐个写 \code{@JsonProperty}，各适合什么情况？
\end{enumerate}
\end{think}
